\section{Asesmen Akhir Komprehensif}

\subsection{Petunjuk Umum}

Asesmen ini dirancang untuk mengukur pencapaian CPMK-1, CPMK-2, CPMK-3, dan CPMK-4 beserta seluruh Sub-CPMK (1.1--1.3, 2.1--2.3, 3.1--3.3, 4.1--4.3) secara menyeluruh \cite{rosen2019}. Instrumen Bagian A--D dapat dipetakan ke komponen Kuis, UTS, UAS, atau proyek akhir sesuai kebijakan RPS. Kerjakan dengan jujur dan mandiri.

\textbf{Alokasi Waktu}:
\begin{itemize}
  \item Bagian A (Pilihan Ganda): 30 menit
  \item Bagian B (Tugas Praktik): 45 menit
  \item Bagian C (Tabel Kebenaran \& Bukti): 60 menit
  \item Bagian D (Pemodelan Praktik Basis Pengetahuan): 45 menit
\end{itemize}

\subsection{Bagian A: Pilihan Ganda (Sub-CPMK 1.1--1.3, 2.1--2.3, 3.1--3.3, 4.3)}

\textbf{Petunjuk}: Pilih satu jawaban yang tepat.

\begin{enumerate}
  \item Manakah yang merupakan proposisi? (Sub-CPMK 1.1)
  \begin{enumerate}
    \item Berapa harga buku ini?
    \item $2 + 2 = 4$
    \item Tolong buka jendela
    \item $x > 5$
  \end{enumerate}
  
  \item Nilai $p \rightarrow q$ ketika $p$ salah dan $q$ benar adalah: (Sub-CPMK 1.2)
  \begin{enumerate}
    \item Benar
    \item Salah
  \end{enumerate}
  
  \item $\neg(p \wedge q)$ ekuivalen dengan: (Sub-CPMK 1.2)
  \begin{enumerate}
    \item $\neg p \wedge \neg q$
    \item $\neg p \vee \neg q$
    \item $p \vee q$
    \item $p \wedge q$
  \end{enumerate}
  
  \item Negasi dari $\forall x \, P(x)$ adalah: (Sub-CPMK 2.2)
  \begin{enumerate}
    \item $\forall x \, \neg P(x)$
    \item $\exists x \, \neg P(x)$
    \item $\neg \exists x \, P(x)$
    \item $\exists x \, P(x)$
  \end{enumerate}
  
  \item Aturan inferensi ``jika $p \rightarrow q$ dan $p$, maka $q$'' disebut: (Sub-CPMK 1.3)
  \begin{enumerate}
    \item Modus Tollens
    \item Modus Ponens
    \item Silogisme
    \item Kontraposisi
  \end{enumerate}
  
  \item Dalam Pemodelan Logika Predikat, fungsi proposisional ber-Arity 2 lazim digunakan untuk merepresentasikan: (Sub-CPMK 4.3)
  \begin{enumerate}
    \item Konstanta atomik
    \item Relasi antar dua entitas individu
    \item Hasil pengujian sintaks
    \item Nilai kebenaran ganda
  \end{enumerate}
  
  \item Manakah yang benar tentang karakteristik fungsi model dalam First-Order Logic? (Sub-CPMK 4.3)
  \begin{enumerate}
    \item Simbol Predikat ditulis dengan huruf kapital awal, diikuti argumen konstanta di dalam kurung
    \item Hanya bisa menggunakan variabel diskrit biner
    \item Tidak dapat menangani operator implikasi
    \item Harus melalui kompilasi interpreter
  \end{enumerate}
  
  \item Ekspresi Boolean $A \wedge B$ direpresentasikan oleh gerbang logika: (Sub-CPMK 3.3)
  \begin{enumerate}
    \item Gerbang AND
    \item Gerbang OR
    \item Gerbang NOT
    \item Gerbang NAND
  \end{enumerate}
  
  \item Dalam verifikasi kebenaran program, Hoare Triple $\{p\}\ S\ \{q\}$ menyatakan: (Sub-CPMK 4.3)
  \begin{enumerate}
    \item $p$ dan $q$ selalu bernilai benar setelah $S$ dijalankan
    \item Jika prekondisi $p$ dipenuhi dan $S$ berterminasi, maka postkondisi $q$ terpenuhi
    \item $S$ pasti berterminasi untuk setiap masukan
    \item Loop invariant sama dengan $q$
  \end{enumerate}
\end{enumerate}

\subsection{Bagian B: Tugas Praktik (Sub-CPMK 2.1--2.3, 3.1, 3.2, 4.1, 4.2, 4.3)}

\textbf{Petunjuk}: Jawab dengan jelas dan lengkap.

\begin{enumerate}
  \item Jelaskan perbedaan tautologi, kontradiksi, dan kontingensi! Berikan contoh masing-masing. (Sub-CPMK 1.2)
  \item Terjemahkan ke logika predikat: ``Semua mahasiswa yang rajin belajar akan lulus.'' Tentukan domain dan predikat yang dipilih. (Sub-CPMK 2.1, 2.2)
  \item Jelaskan metode bukti langsung, kontraposisi, dan kontradiksi. Dalam konteks apa masing-masing metode biasanya dipilih? (Sub-CPMK 4.1)
  \item Jelaskan prinsip induksi matematika (basis dan langkah induktif). Berikan contoh pernyataan yang dapat dibuktikan dengan induksi. (Sub-CPMK 4.2)
  \item Jelaskan secara konseptual tiga tumpuan pilar Logika Predikat dalam dunia industri kecerdasan buatan, terutama mengenai Entitas, Kuantor/Aturan, dan Inferensi Resolusi. (Sub-CPMK 2.1--2.3, 4.3)
  \item Jelaskan perbedaan kebenaran parsial dan kebenaran total dalam verifikasi program. Mengapa loop invariant saja tidak cukup untuk menjamin kebenaran total? Berikan contoh program yang memenuhi kebenaran parsial tetapi tidak berterminasi. (Sub-CPMK 4.3)
\end{enumerate}

\subsection{Bagian C: Tabel Kebenaran dan Bukti (Sub-CPMK 1.2, 1.3, 4.1, 4.2)}

\textbf{Petunjuk}: Kerjakan secara sistematis. Tulis langkah-langkah bukti dengan jelas \cite{forallx,hammack2018}.

\begin{enumerate}
  \item Buat tabel kebenaran untuk $(p \rightarrow q) \leftrightarrow (\neg p \vee q)$. Apakah ekuivalensi ini tautologi? (Sub-CPMK 1.2, 1.3)
  \item Buktikan dengan induksi: $1 + 2 + \cdots + n = n(n+1)/2$ untuk $n \geq 1$. (Sub-CPMK 4.2)
  \item Buktikan: ``Jika $n^2$ genap maka $n$ genap'' dengan metode kontraposisi. (Sub-CPMK 4.1)
\end{enumerate}

\subsection{Bagian D: Pemodelan Praktik Basis Pengetahuan (Sub-CPMK 4.3)}

\textbf{Petunjuk}: Kerjakan pemodelan skenario dunia nyata di bawah ini menggunakan gaya notasi Aljabar Logika Predikat Orde Pertama (FOL).

\textbf{Studi Kasus: Pemodelan Sistem Pendeteksian Spam}

Buat konstruksi formal struktur logika untuk mendeteksi apakah sebuah \textit{Email} dikategorikan Spam berdasarkan requirements berikut:

\textbf{Requirements}:
\begin{enumerate}
  \item \textbf{Fakta Dasar Email}: Definisikan minimal 5 fakta relasi \texttt{BerisiKata(EmailID, KataKunci)} untuk menyatakan sebuah email memuat vocabulary tertentu. Contoh: \texttt{BerisiKata(Email01, "Diskon")}.
  \item \textbf{Fakta Konteks Pengirim}: Definisikan 3 fakta sifat unary (Arity 1) mengenai pengirimnya mencurigakan, misalnya: \texttt{UnknownSender(Email01)}.
  \item \textbf{Aturan (Rule of Inference)}: Modifikasi sebuah aturan \texttt{IsSpam(x)} berkuantor universal ($\forall$) yang menyatakan: ``Bila Email x mengandung kata `Diskon' dan ia datang dari UnknownSender, maka x adalah Spam''. Tuliskan persamaan logikanya secara tepat.
  \item \textbf{Simulasi Deduksi}: Misalkan \texttt{Email99} punya 2 atom kebenaran: berisi kata `Diskon' dan dikirim oleh UnknownSender. Buktikan secara modus ponens via substitusi formal bahwa \texttt{IsSpam(Email99)} terpenuhi.
\end{enumerate}

\textbf{Kriteria Penilaian}: Kebenaran sintaks fungsi (kapitalisasi, nesting fungsi), kelengkapan atom kebenaran, validitas penurunan rantai deduktif, serta kesesuaian dengan teori Logika Predikat Orde Pertama.
